%! Characteristics of Radical Functions — Unit 5, Lesson 5.0 — Student Packet Cover Sheet
\documentclass[10pt]{article}
\usepackage{algebra2-article}
\usepackage{algebra2-boxes}
\usepackage{ltablex}
\keepXColumns

\begin{document}

% Full-bleed forest banner
\begin{tikzpicture}[remember picture, overlay]
  \fill[forest] (current page.north west)
    rectangle ([yshift=-0.9in]current page.north east);
\end{tikzpicture}
\vspace{-0.8in}

\begin{center}
  {\color{white}\textbf{\LARGE Algebra 2: Shepherd}}\\[4pt]
  {\color{forestmid}\large\textbf{Unit 5 \quad Radical Functions}}\\[6pt]
  {\color{white}\textbf{\large Lesson 5.0 \quad Characteristics of Radical Functions}}
\end{center}

\vspace{0.22in}
\namedateperiod
\vspace{0.18in}

\begin{learningtargetbox}
\small
\begin{enumerate}[leftmargin=1.5em, itemsep=5pt]
  \item \ldots{} find a radical function's \textbf{domain} from its \textbf{index} and \textbf{radicand}, and explain why an \textbf{even} index forces a restricted domain while an \textbf{odd} index does not.
  \item \ldots{} read the \textbf{domain}, \textbf{range}, \textbf{endpoint}, \textbf{intercepts}, \textbf{increasing/decreasing} and \textbf{positive/negative} intervals, \textbf{absolute extrema}, and \textbf{end behavior} off a square root or cube root graph.
  \item \ldots{} explain the \textbf{inverse relationship} between $y=\sqrt{x}$ and $y=x^2$ (restricted) and between $y=\sqrt[3]{x}$ and $y=x^3$, and \textbf{compare and contrast} the two radical families.
\end{enumerate}
\end{learningtargetbox}

\vspace{0.14in}

\begin{tocbox}
\small
\renewcommand{\arraystretch}{1.5}
\begin{tabularx}{\linewidth}{@{} c l X r @{}}
\rowcolor{forestbg}
\textbf{\#} & \textbf{Component} & \textbf{Description} & \textbf{Score} \\
\midrule
1 & Warm-Up        & Spiral review: which roots exist, solving \emph{radicand} $\ge 0$, and a pair of swapped tables & \blank{1.2cm} \\
2 & Guided Notes   & Reading every characteristic off the square root and cube root parents            & \blank{1.2cm} \\
3 & Group Activity & \emph{Reading Radical Graphs} --- feature tables, domains from equations, and a falling-object model, by tier & \blank{1.2cm} \\
4 & Exit Ticket    & Quick check: features from a graph, plus an SOL-style domain item & \blank{1.2cm} \\
5 & Homework       & Domains from equations, two contrasting graphs, and a skid-mark model    & \blank{1.2cm} \\
\midrule
  & \multicolumn{2}{r}{\textbf{Total}} & \blank{1.2cm} \\
\end{tabularx}
\end{tocbox}

\vspace{0.10in}

\begin{remindbox}
\small A \textbf{radical function} puts the variable under a radical, $f(x)=\sqrt[n]{\phantom{x}\,}$,
and the \textbf{index} $n$ decides everything. An \textbf{even} index (a plain $\sqrt{\phantom{x}}$)
refuses a negative \textbf{radicand}, because no real number squares to a negative --- so the domain
is whatever solves \emph{radicand} $\ge 0$, and the graph begins at an \textbf{endpoint} and travels
in \emph{one} direction only. That endpoint is always an \textbf{absolute extremum}: a minimum if the
arm climbs away from it, a maximum if the arm falls. An \textbf{odd} index (like $\sqrt[3]{\phantom{x}}$)
refuses nothing, because cubing keeps the sign --- so the domain and range are \emph{all reals}, there
is no endpoint, no extrema, and the graph is symmetric about the \textbf{origin} with two arms.
Both parents come from somewhere you have already been: $y=\sqrt{x}$ is $y=x^2$ (restricted to
$x\ge 0$) reflected over the line $y=x$, and $y=\sqrt[3]{x}$ is all of $y=x^3$ reflected the same way.
That is what an \textbf{inverse} is --- inputs and outputs traded places --- and it is the reason a
square root only accepts the numbers a square can produce.
\end{remindbox}

\end{document}
